\chapter{Preface}

\section*{Who this book is for}

You want to make your own artificial life. Maybe you watched a Lenia
animation and wondered how it works, or saw a Boids flock and wondered
what changes if you alter the rules, or read about Tierra and wondered
whether you could build something that does what it did but does not
stop.

You can read a probability density without flinching, you have seen at
least one dynamical system in motion, and you have written enough
Python to follow a code listing. That is the technical baseline; you
do not need anything beyond it. The rest is curiosity.

This book offers you a vocabulary --- five words: \emph{substrate},
\emph{variation}, \emph{viability}, \emph{topology}, \emph{observer}
--- and shows you how to use them to read, design, and compare systems.
The vocabulary will not write your simulator for you. It will tell you
what choices you are making when you write one, and what choices the
people whose work you are reading have made. That turns out to be the
harder part.

\section*{What this book is not}

This is not a survey of every artificial-life system. There are
excellent surveys --- Bedau and Cleland's \textit{The Nature of Life},
Langton's \textit{Artificial Life} proceedings, and Aguilar et al.'s
recent reviews --- and we will point you toward them. We mention only
the systems that illuminate something specific about the vocabulary.

This is also not a philosophy-of-life book, though one chapter does
take philosophy of life seriously. The five-slot vocabulary is
descriptive: it lets you read what people mean by \emph{life},
\emph{emergence}, and \emph{novelty} without forcing you to take sides.
Several philosophical frames will reappear later as concrete technical
formalisms, and we will point that out, but the book does not tell you
which is right.

Finally, this is not a manual for building one particular
artificial-life system. There is a code companion --- a typed JAX
scaffold --- that goes alongside it, and the worked examples in each
chapter are runnable. But the book is after a way of reading and
designing systems, not a recipe for any one of them.

\section*{The central thesis}

Artificial-life systems can be read as a five-tuple: a \emph{substrate}
of states, a \emph{variation} operator that generates new candidates, a
\emph{viability} filter that decides what persists, an
\emph{interaction topology} governing who-affects-whom, and an
\emph{observer} that measures what we can actually see. We will call
this the five-slot lens.

The lens is descriptive, not prescriptive. It does not pronounce on
what counts as alive, what counts as emergent, or what counts as novel.
It gives you a structured place to write each candidate answer, so you
can compare frames that would otherwise pass each other in the night.

This stance --- \emph{vocabulary, not theory} --- is the choice that
organises everything else in the book. We will not pretend the four
formalisms of viability collapse into one (there is an open conjecture
about whether they do). We will not pretend effective information has
crowned itself the definition of emergence (it is a useful measure
under a stated convention). We will not pretend $\Omega$ solves
open-endedness (it is one proposal with stated weaknesses).

\section*{A map of what is coming}

The book has three movements.

\textbf{Phenomena and puzzles (Chapters 1--3).} Chapter 1 tours a
handful of artificial-life systems --- Game of Life, Boids, Lenia,
Tierra, evolutionary computation --- without committing to a vocabulary
for them. Chapter 2 takes the philosophy of life seriously: Aristotle
through Schr\"odinger through Friston, with the contemporary frames we
will meet again later. Chapter 3 surveys what people mean by
\emph{emergence} and \emph{novelty}, and where the words conflict.

\textbf{The vocabulary (Chapters 4--7).} Chapter 4 introduces the
five-slot lens and re-reads the systems from Chapter 1 in it. Chapters
5--7 develop each slot --- substrate and variation, then viability
(with the four formalisms), then topology and observers --- with the
running example of Lenia carried through.

\textbf{Composition, measurement, and frontiers (Chapters 8--9).}
Chapter 8 shows how to build bigger systems out of smaller ones (the
coupled-Lenia worked example) and how to measure emergence and novelty
quantitatively. Chapter 9 surveys what is still open: conjectures,
frames the lens has not yet absorbed, and the lens applied beyond
engineered systems.

A math primer at the back of the book introduces the four pieces of
formal machinery the chapters lean on: order theory, measure theory,
information theory, and a light dose of category theory. The primer
is self-contained and meant to be readable cold; flip to it when a
chapter writes down a symbol you would like to read more slowly.

A reader who wants the formal definitions before the motivation can
start at Chapter 4 and reach back to Chapters 1--3 as needed. A
reader who wants the math foundations first can start at the primer.
The book is written for linear reading, but you will not be punished
for jumping.

\section*{Conventions}

\textit{Notation.} The five slots take the symbols $X$ (substrate),
$V$ (variation), $F$ (viability), $T$ (topology), $O$ (observer). State
spaces are typeset in plain capitals; distributions over states are
written $\mathcal{P}(X)$, and signed measures are written
$\mathcal{M}_{+}(X)$. When a Greek letter appears in the formal
literature --- $\Phi$ for substrate dynamics, $\mu$ for the initial
measure, $\kappa$ for substrate coupling, $\rho$ for variation-channel
composition, $\Omega$ for open-endedness --- we use it, with the
plain-English name on first appearance.

\textit{Code companion.} Each chapter points to a notebook in the JAX
scaffold that accompanies the book. The scaffold is typed,
jit-friendly, and small; it does not ship concrete systems, only the
harness. Running the notebooks is optional but recommended for at
least one system per chapter.

\textit{End-of-chapter prompts.} Rather than formal exercise sets, each
chapter ends with a handful of ``to think about'' prompts --- a mix of
conceptual, computational, and open-ended. They are the closest thing
to homework the book contains.

\textit{Citations.} Inline and selective rather than exhaustive. The
``further reading'' at the back is curated and annotated.

\section*{Thanks}

[Placeholder.]
